\chapter*{Conclusion}
\addcontentsline{toc}{chapter}{Conclusion}

Ce mémoire avait pour objectif de déterminer dans quelle mesure la modélisation thématique et la classification automatique permettent d'identifier des marqueurs du naturalisme sans confondre ce mouvement avec la signature de Zola, les sujets de ses romans ou les effets de constitution du corpus. Les trois chapitres ont abordé cette question de manière progressive : les deux premiers ont dégagé des régularités internes à l'œuvre zolienne, tandis que le troisième a cherché à savoir si une partie de ce signal se retrouvait chez d'autres auteurs.

Le premier chapitre a montré que la NMF et K-means pouvaient organiser les trente et un romans de Zola en ensembles lexicaux interprétables. Le travail, la mine, l'argent, la religion, la famille, la guerre, le commerce et les espaces domestiques ou urbains réapparaissent dans les deux analyses. Cette convergence reste toutefois exploratoire : les méthodes utilisent les mêmes segments et la même matrice TF--IDF, tandis que les 29 clusters de K-means ont été retenus en partie pour conserver une granularité comparable aux 29 topics de la NMF. Elles décrivent donc des régularités internes à Zola, sans se valider indépendamment l'une l'autre.

Le deuxième chapitre a retrouvé plusieurs de ces domaines avec BERTopic, malgré une segmentation et une représentation différentes. Les correspondances sont les plus nettes pour les vocabulaires spécialisés de la finance, de la mine, du chemin de fer, de la guerre ou de l'alimentation populaire. L'analyse chronologique fait aussi apparaître des recompositions : les relations affectives sont très présentes dans plusieurs œuvres de jeunesse, tandis que les institutions religieuses, politiques et judiciaires occupent une place plus importante dans les cycles tardifs. Ces observations ne décrivent cependant pas une évolution continue de l'écriture. Certains sommets sont principalement portés par un seul roman, et plus des trois quarts des segments ont reçu leur topic après une réaffectation aux noyaux initialement découverts. Il me semble donc plus juste de considérer les 17 topics comme une cartographie exploratoire que comme des catégories naturelles du corpus.

Le troisième chapitre a déplacé l'analyse vers un corpus comparatif. Le modèle entraîné avec Zola comme unique auteur positif classe correctement 18 œuvres sur 25 dans la condition avec noms propres et 17 sur 25 après leur suppression automatique. Les AUC moyennes par auteur restent respectivement de $0{,}859$ et $0{,}836$, et l'ordre des scores des deux conditions est très proche, avec une corrélation de rang de $0{,}972$. Dans cet échantillon, le signal ne disparaît donc pas lorsque les noms propres détectés sont retirés. Je ne peux cependant pas affirmer qu'il s'agit d'un style naturaliste autonome : le rappel des œuvres naturalistes reste faible, l'effet varie selon les auteurs et les contrôles chronologiques sont trop réduits pour exclure un effet de période.

Au regard de ces résultats, il me semble que la réponse à la problématique est partiellement positive. Les méthodes employées permettent de repérer des marqueurs candidats, d'en observer la réapparition à travers plusieurs chaînes de traitement et de tester leur circulation hors de l'œuvre de Zola. Elles ne permettent pas encore de séparer complètement contenu thématique, signature d'auteur et artefacts de corpus. Le résultat le plus solide n'est donc pas la découverte d'une définition automatique du naturalisme, mais la mise en évidence d'une proximité textuelle mesurable entre Zola et certaines œuvres naturalistes d'autres auteurs.

Cette conclusion doit rester liée au caractère exploratoire du travail. Le premier corpus ne comprend qu'un auteur, le modèle de classification n'apprend la classe naturaliste qu'à partir de Zola, et le test comparatif ne porte que sur trois auteurs. Les catégories littéraires de certaines œuvres peuvent également être discutées. Enfin, plusieurs choix de segmentation, de nettoyage, de nombre de groupes et de seuil ont été guidés en partie par l'interprétabilité. Ces limites ne rendent pas les résultats inutiles ; elles définissent plutôt ce qu'ils permettent réellement d'avancer et ce qui reste à vérifier.

Une première suite consisterait à constituer un nouveau test réellement aveugle, plus large et mieux équilibré selon les auteurs et les périodes. L'apprentissage pourrait intégrer plusieurs écrivains naturalistes, puis laisser successivement un auteur entier hors du modèle. Un témoin reposant uniquement sur l'année de publication aiderait également à mesurer ce que le texte apporte au-delà de la chronologie. Je pense que cette étape serait prioritaire avant de tester des modèles plus complexes.

Il serait ensuite utile de confronter séparément le contenu et la forme. Les topics dégagés dans les deux premiers chapitres pourraient être recherchés dans le corpus comparatif, tandis qu'une autre expérience utiliserait des mots-outils, la ponctuation, les n-grammes de caractères ou les catégories grammaticales. La suppression des noms propres pourrait aussi être évaluée manuellement et comparée à leur remplacement par des catégories génériques. Ces prolongements ne garantissent pas qu'un style naturaliste stable puisse être isolé, mais ils permettraient de préciser progressivement la part respective des thèmes, de l'époque et des habitudes d'écriture dans les résultats observés.
